\chapter{Theoretical Framework}\label{theoretical-framework}

This section provides the mathematical vocabulary and analytical tools
necessary in answering the three research questions guiding this
research. Specifically, it establishes: (1) how to formally specify
IBC\textquotesingle s atomic actions and processes, (2) how to extract
observable behaviours including traces, deadlocks, and branching
structures from such specifications, (3) how to operationally determine
whether asset and data behaviours within IBC are unified, stacked, or
conflicting using behavioural equivalence.

\section{Systems and System
Behaviour}\label{systems-and-system-behaviour}

A \textbf{system} is formally composed of two major aspects: process and
data. In this framework, processes serve as control mechanisms that runs
and manipulates data. Meanwhile, data remains static and passively
comprises the system \citep{fokkink2007introduction}. The composition of
these elements formsa structure that can perform operations according to
specified rules \citep{plotkin2004structural}.

The \textbf{behaviour} of a system refers to the sum of its ability to
interact or communicate. A behaviour is \textbf{observable} because it
can be seen from its outside operations; particularly, the interaction
of the inputs it accepts and the outputs it produces
\citep{milner1989communication}. Every elementary step or operation that
happens in the system is called \textbf{transition} that changes the
system from its current configuration to another state
\citep{plotkin2004structural}.

A \textbf{Blockchain} is an example of a system that can be treated as a
transaction-based \textbf{state machine} \citep{wood2021ethereum}. In
this view, the blockchain is not just a storage of information but a
dynamic mechanism where transactions serve as inputs that result to
transitions of the distributed ledger\textquotesingle s state. Every
succesful transaction creates new state that is formally accepted by the
whole network. The usage of \textbf{formal modelling} is critical in
this topic in order to give sufficient mathematical evidence for the
security and stability of interoperability systems. Formal modelling
removes ambiguity and ensures the correctness in comparing asset-bsed
and data-based mechanisms under one theoretical framework
\citep{Duan2018} \citep{fokkink2007introduction}.

\section{Interoperability as System
Composition}\label{interoperability-as-system-composition}

In process theory, interoperability can be formally described as a
composition of two or more systems that work together to achieve a
particular goal \citep{lalli2024}. Under the framework of Process
Algebra, this communication is symbolised by the parallel composition
\[ (P || Q) \], wherein every blockchain network is treated as separate
process that runs concurrently and independently from the other networks
but with the ability to communicate to each other.

The communication inside this composition is run by the mechanism of
\textbf{synchronisation} \citep{milner1989communication}. By means of a
formal \textbf{communication function}, the individual actions of every
blokchain gets to communicate. For example, the action that delivers the
message \((send)\) from a system needs to match the action of receiving
\((read)\) at the other system in order to complete a succefull
communication.

The behaviour of the formed composition is treated as observable from
the perspective of an outsider viewer \citep{milner1989communication}.
Using the \textbf{abstraction operator} the complex details of the
internal communication between blockchains are hidden or made as
\emph{silent steps (\(\tau\)),} so that what only remains are the
external input-output relation of the whole inter-operational system
\citep{fokkink2007introduction}. The resulting traces or series
observable actions serve as the evidence if the operation of the
composed systems remain consistent to the original semantic
specification.

This framework of formal composition gives way to formally analyse and
compare the two ways of blockchain communications: the asset-based and
data-based interoperability.

\section{Foundations of Process
Algebra}\label{foundations-of-process-algebra}

\textbf{Process Algebra} serve as the main theoretical instrument of
this research particularly the principles of communication, composition,
and abstraction. The research will use it to formally analyse and
manipulate operations of blockchains. This is illustrated as the study
of concurrent systems or communicating systems inside an axiomatic and
algebraic framework \citep{bergstra1982fixed}
\citep{fokkink2007introduction}. Instead of just using visual diagrams,
processes are embodies as formal \emph{process terms} using a
mathematical language that allows linear notation and mechanised
manipulation \citep{fokkink2007introduction}.

\subsection{Atomic Actions and
Alphabet}\label{atomic-actions-and-alphabet}

System behaviour is constructed from a formally defined alphabet of
atomic actions. Let: \[ D: \text{set of data} \]
\[ P: \text{set of communication ports} \] Atomic actions are defined
as: \[ Send: s_p(d) \] \[ Read: r_p(d) \] where \(d \in D\) and
\(p \in P\) \citep{Vaandrager_1990}.

\subsection{Communication and
Handshaking}\label{communication-and-handshaking}

Communication between processes is governed by a communication function:
\[ \gamma(s_p(d), r_p(d)) = c_p(d) \] This formalises the handshaking
paradigm where communication happens only through synchronisation
between matching send and read actions. Isolated actions are disallowed
to ensure correctness in system interaction \citep{Vaandrager_1990}.

\subsection{Process Operators}\label{process-operators}

In this framework, every system\textquotesingle s process is composed of
atomic actions (actions that has no duration) combined using the
following operators: \textbf{Sequential Composition \(( x \cdot y )\):}
\(x\) is perfomed first before \(y\). \textbf{Alternative Composition
\(( x + y )\):} symbolises the nondeterministic choice between \(x\) and
\(y\) operations. \textbf{Parallel Composition \(( x || y )\):}
symbolises the simultaneous execution of actions \(x\) and \(y\).

A fundamental feature of process algebra is its
\textbf{compositionality} wherein a large system can be treated as
combined smaller and simpler parts \citep{fokkink2007introduction}.
Using the \textbf{parallel operator} ( \(||\) ), individual blockchain
network or protocol components can be modelled separately and then
regouped to see the whole behaviour of the inter-operational system
\citep{fokkink2007introduction}. This modular technique allows the
analysis of the connection between asset-based and data-based protocols
as a series of concurrent processes \citep{fokkink2007introduction}.

\subsection{Encapsulation}\label{encapsulation}

To control system visibility: \textbf{Encapsulation (\(\partial_H\))}
blocks unmatched internal actions. It is used in order to ensure these
internal actions do not happen alone but as a part if a synchronized
communciation that causes transition to the state of the whole
composition.

\subsection{Recursion and Abstraction}\label{recursion-and-abstraction}

Recursion allows specification of infinite or repetitive behaviours
using fixed-point operators. For blockchain protocols, recursion models
continuous channel operation:

\[
Channel = Open \cdot (Send \cdot Receive + Timeout) \cdot Channel
\]

The \textbf{silent action (\(\tau\))} represents internal, unobservable
steps. Relayer forwarding, proof verification, and internal state
updates are modelled as \(\tau\) actions. That is, they affect the
system\textquotesingle s internal state but are not visible to an
external observer.

The \textbf{abstraction operator \(\tau_I\)} renames all actions in set
\(I\) to \(\tau\), hiding implementation details. For IBC:

\[
\tau_{\{forward\_packet, verify\_proof\}}(P_{IBC})
\]

makes only packet sends, receives, acknowledgements, and timeouts
observable. Meanwhile, internal relayer behaviour becomes silent steps.

Key properties of \(\tau\):

\begin{itemize}
    \item \(\tau \cdot P = P\) (silent steps can be removed without changing observable behaviour)
    \item \(a \cdot (\tau \cdot (b + c))\) is observationally equivalent to \(a \cdot (b + c)\)
\end{itemize}

\subsection{Operation Semantics}\label{operation-semantics}

The meaning of process expressions is formally defined using
\textbf{Structural Operational Semantics (SOS)}. Transition rules
specify how processes evolve. This ensures that every behaviour is
derivable and mathematically well-defined \citep{Vaandrager_1990}.

\section{Labelled Transition Systems as Semantic
Model}\label{labelled-transition-systems-as-semantic-model}

Labelled Transition System (LTS) is a fundamental mathematical structure
used to model computations and dynamic behaviour of systems
\citep{vanBenthemBergstra1995}. A Labelled Transition System (LTS) is
formally defined as a tuple \((S, I, Act, \rightarrow, \checkmark)\)
where:

\begin{itemize}
    \item \(S\) is a set of states
    \item \(I \in S\) is the initial state
    \item \(Act\) is a set of action labels
    \item \(\rightarrow \subseteq S \times Act \times S\) is the transition relation
    \item \(\checkmark \subseteq S\) is the set of \textbf{successfully terminated states}
\end{itemize}

A \textbf{deadlock state} is a state with no outgoing transitions
(\(\nexists a, s'\) such that \(s \xrightarrow{a} s'\)) AND
\(s \notin \checkmark\). Deadlock represents protocol failure wherein a
packet can neither be delivered nor timed out.

For IBC, deadlock occurs when:

\begin{itemize}
    \item A handshake remains incomplete (OpenInit sent, no OpenTry received)
    \item A packet is sent but neither acknowledged nor timed out
\end{itemize}

The \textbf{observable behaviour} of the system is composed of the
\textbf{traces}. They are the series of actions from the starting state
\citep{Gorrieri_2017}. This formalism allows the abstraction from the
inside implementation to the outside communication
\citep{plotkin2004structural}. These traces serve as the guidelines to
analyse whether a protocol follows its intended semantic
interoperability.

It is important to notice when two systems with differing structure
representation or graph can share the same behaviour
\citep{baeten1995concrete}. This concept is called \textbf{behavioural
equivalence} wherein the states of different LTSs are treated as equal
if they exhibit the same pattern of transitions such as what the theory
of bisimulation describes \citep{baeten1995concrete}
\citep{Gorrieri_2017}. This abstraction allows the inquiry of comparison
of protocols for asset-based and data-based interoperability in
blockchains without limiting particular details of their internal
implementation \citep{plotkin2004structural}.

Nevertheless, the simple usage of state transition graphs can lead to
the problem of state explosion that makes formal verification difficult
\citep{vaandrager1986verification}. For this reason, there is a need to
a more expressive framework that allows algebraic manipulation and
composition of processes.

\section{Behavioural Equivalences}\label{behavioural-equivalences}

In order to identify whether an implementation of a protocol matches to
its expected behaviour, the concept of \textbf{behavioural equivalence}
is used; particularly, the \textbf{branching bisimulation} (
\(\leftrightarrow_{rb}\) ) \citep{Fokkink2004}
\citep{fokkink2007introduction}. Branching simulation is needed because
it maintains the branching structure of processes that is crucial in
identifying deadlocks which can result to loss of value or message
\citep{Fokkink2004}. In order to prove this bisimulation between two
different ledgers, \textbf{Cones and Foci methodology} must be used
where every state in an implementation is mapped to a particular focus
point at a specification that shows similar external behaviour
\citep{Fokkink2004}.

\subsection{Trace Equivalence (Partial and
Completed)}\label{trace-equivalence-partial-and-completed}

A \textbf{partial trace} is any finite sequence of observable actions
from the initial state. On the otherhand, a \textbf{completed trace} is
a sequence ending in a deadlock or termination state.

Two processes \(P\) and \(Q\) are \textbf{trace-equivalent}
(\(P \approx_{tr} Q\)) if \(Traces(P) = Traces(Q)\).

\textbf{Limitation}: Trace equivalence collapses branching structure.
The processes \(a.(b + c)\) (choose between b or c after a) and
\(a.b + a.c\) (choose between doing a then b or a then c) have identical
trace sets \(\{ab, ac\}\) but exhibit different branching behaviour at
the moment of choice.

For interoperability, trace equivalence cannot distinguish a protocol
that offers a choice between acknowledgement and timeout
\((ack + timeout)\) from one that forces acknowledgement then timeout
\((ack \cdot timeout)\).

\subsection{Bisimulation (Strong and
Weak)}\label{bisimulation-strong-and-weak}

\textbf{Strong bisimulation} (\(\leftrightarrow\)) requires that every
action of \(P\) is matched by exactly the same action of \(Q\),
including \(\tau\). A relation \(R\) is a strong bisimulation if
whenever \((P,Q) \in R\):

\begin{itemize}
    \item If \(P \xrightarrow{a} P'\), then \(\exists Q'\) such that \(Q \xrightarrow{a} Q'\) and \((P',Q') \in R\)
    \item If \(Q \xrightarrow{a} Q'\), then \(\exists P'\) such that \(P \xrightarrow{a} P'\) and \((P',Q') \in R\)
\end{itemize}

\textbf{Weak bisimulation} (\(\approx\)) abstracts over \(\tau\) steps.
Define \(\Rightarrow\) as zero or more \(\tau\) steps. Then
\(P \xrightarrow{a} \Rightarrow Q\) means
\(P \Rightarrow \xrightarrow{a} \Rightarrow Q\). Weak bisimulation
requires matching actions after absorbing \(\tau\) steps.

For IBC, weak bisimulation is appropriate because relayer forwarding
(\(\tau\)) should not affect observable correctness.

\subsection{Branching Bisimulation and Cones and
Foci}\label{branching-bisimulation-and-cones-and-foci}

\textbf{Branching bisimulation} (\(\leftrightarrow_{rb}\)) preserves the
branching structure of processes, unlike trace equivalence. It is
crucial for identifying deadlocks which can result in loss of value or
messages \cite{Fokkink2004}.

To prove bisimulation between two processes (e.g., asset and data
specifications), this research adopts the
\textbf{Cones and Foci methodology} \cite{Fokkink2004}:

\begin{itemize}
    \item A \textbf{cone} is the set of all states reachable from a given state
    \item A \textbf{focus} is a state in the specification that the implementation maps to
    \item Construct a state mapping function \(\phi\) from implementation states to specification states
    \item Prove each \(\tau\)-step either preserves \(\phi\) or moves closer to a focus point
\end{itemize}

This methodology establishes branching bisimulation between
IBC\textquotesingle s concrete behaviour and its abstract specification.

\section{From Process Algebra to LTS: Operational
Semantics}\label{from-process-algebra-to-lts-operational-semantics}

Structural Operational Semantics (SOS) defines how process terms
transition. Each operator has inference rules \citep{Vaandrager_1990}.

\subsection{Atomic Action}\label{atomic-action}

\[
\frac{}{a \xrightarrow{a} \checkmark}
\]

\subsection{\texorpdfstring{Sequential Composition
(\(P \cdot Q\))}{Sequential Composition (P \textbackslash cdot Q)}}\label{sequential-composition-p-cdot-q}

\[
\frac{P \xrightarrow{a} P'}{P \cdot Q \xrightarrow{a} P' \cdot Q} \quad
\frac{P \xrightarrow{a} \checkmark}{P \cdot Q \xrightarrow{a} Q}
\]

\subsection{\texorpdfstring{Alternative Composition
(\(P + Q\))}{Alternative Composition (P + Q)}}\label{alternative-composition-p-q}

\[
\frac{P \xrightarrow{a} P'}{P + Q \xrightarrow{a} P'} \quad
\frac{Q \xrightarrow{a} Q'}{P + Q \xrightarrow{a} Q'}
\]

\subsection{\texorpdfstring{Parallel Composition
(\(P \parallel Q\))}{Parallel Composition (P \textbackslash parallel Q)}}\label{parallel-composition-p-parallel-q}

\[
\frac{P \xrightarrow{a} P'}{P \parallel Q \xrightarrow{a} P' \parallel Q} \quad
\frac{Q \xrightarrow{a} Q'}{P \parallel Q \xrightarrow{a} P \parallel Q'}
\]

\subsection{Communication
(Synchronisation)}\label{communication-synchronisation}

\[
\frac{P \xrightarrow{s_p(d)} P' \quad Q \xrightarrow{r_p(d)} Q'}{P \parallel Q \xrightarrow{c_p(d)} P' \parallel Q'}
\] where \(\gamma(s_p(d), r_p(d)) = c_p(d)\).

\subsection{\texorpdfstring{Encapsulation
(\(\partial_H(P)\))}{Encapsulation (\textbackslash partial\_H(P))}}\label{encapsulation-partial_hp}

\[
\frac{P \xrightarrow{a} P' \quad a \notin H}{\partial_H(P) \xrightarrow{a} \partial_H(P')}
\]

\subsection{\texorpdfstring{Abstraction
(\(\tau_I(P)\))}{Abstraction (\textbackslash tau\_I(P))}}\label{abstraction-tau_ip}

\[
\frac{P \xrightarrow{a} P' \quad a \notin I}{\tau_I(P) \xrightarrow{a} \tau_I(P')} \quad
\frac{P \xrightarrow{a} P' \quad a \in I}{\tau_I(P) \xrightarrow{\tau} \tau_I(P')}
\]

These rules provide the formal basis for extracting LTS models from
process algebra specifications of IBC.

\section{Application to Blockchain
Interoperability}\label{application-to-blockchain-interoperability}

Process algebra has been successfully applied to blockchain systems,
establishing precedent for this research.

\begin{itemize}
\item
  The \(\natural\)-calculus (nu-calculus) is a process calculus embedded
  in Haskell and formalised in Isabelle/HOL, has been used to verify the
  Ouroboros consensus protocol for Cardano. Key innovation: using two
  stacked LTSs to handle data transfer and channel scope modularly
  \citep{jeltsh2019}.
\item
  CCS (Calculus of Communicating Systems) and weak bisimulation has been
  used to verify optimising transformations on smart contracts written
  in Cardano\textquotesingle s Marlowe eDSL. They proved that
  transformed contracts preserve observable behaviour
  (\(P_{original} \approx P_{transformed}\)) \citep{mcilwaine2025}.
\end{itemize}

\textbf{Why Process Algebra for IBC Specifically}:

\begin{enumerate}
    \item \textbf{Concurrent channels}: IBC allows multiple channels between same chains. Parallel operator (\(\parallel\)) models this naturally.
    \item \textbf{Branching choices}: IBC's handshake includes choice points (OpenInit → OpenTry OR timeout). Alternative composition (\(+\)) preserves branching structure.
    \item \textbf{Silent abstraction}: Relayer behaviour is internal. \(\tau\) operator and weak bisimulation allow abstraction while preserving observable correctness.
    \item \textbf{Compositional analysis}: IBC separates TAO from application layer. Process algebra allows verifying assets and data separately, then composing them to detect conflicts.
    \item \textbf{Behavioural equivalence}: The central question (unified/stacked/conflicting) requires weak bisimulation as the mathematical criterion.
\end{enumerate}

\section{Framework Synthesis}\label{framework-synthesis}
